\subsection{A quadratic-field reduction for two signed quotients}
\label{subsec:quadratic-signed-reduction}

\begin{proposition}\label{prop:quadratic-signed-reduction}
For a hypothetical categorification of $R7\_4\_1212$ or
$R7\_4\_1213$, choose positive-second-indicator lifts in its
sphericalization and let $d_i$ be their dimensions. Then
\[
 K=\mathbb Q(d_0,\ldots,d_6)=\mathbb Q(\sqrt5),\qquad
 \delta=\sum_i d_i^2\in
 \{180,360,270\pm90\sqrt5,450\pm90\sqrt5\}.
\]
\end{proposition}
\begin{proof}
The two source Casimir polynomials are respectively
\[
\begin{split}
 &(z-90)(z-20)^2(z-4)(z-2)(z^2-900z+6480),\\
 &(z-90)(z-12)(z-10)(z-6)(z-2)(z^2-900z+6480).
\end{split}
\]
Their traces are $B=1036,1020$, and the least common multiple of
the rational formal codegrees is $180$. As in
Proposition~\ref{prop:relative-trace}, the signed multiplication
matrices are commuting symmetric integer matrices. The formal-codegree
field and divisibility theorems for the sphericalization
\cite[Theorem~2.13 and Corollaries~2.14--2.15]{OstrikRank3}
therefore give $\mathbb Q(\sqrt5)\subseteq K$ and
$\beta=\delta/180\in\mathcal O_K$. Every embedding of $K$ is real,
and $\beta$ is totally positive.

The nontrivial invertible $g$ fixes all five noninvertibles in the
source. The square-channel congruence gives $T_{ii}^g=1$, and
Frobenius reciprocity gives $gX_i=X_i$ in the signed quotient. Thus
$(1,-1,0,\ldots,0)$ is a character of codegree two, distinct from
every conjugate of the dimension character. If $e=[K:\mathbb Q]$,
then $e$ is even, $e\leq6$, and
\[
 180\operatorname{Tr}_{K/\mathbb Q}(\beta)+2\leq B.
\]
Total positivity and the arithmetic--geometric mean inequality give
$\operatorname{Tr}_{K/\mathbb Q}(\beta)\geq e$. Hence $e\ne6$.

Suppose $e=4$ and put $r=[\mathbb Q(\beta):\mathbb Q]\mid4$.
If $r=4$, its minimal polynomial is
$z^4-az^3+bz^2-cz+d$, with positive integer coefficients and
\[
 a\leq5,\qquad b\leq\lfloor3a^2/8\rfloor,\qquad
 c\leq\lfloor a^3/16\rfloor,\qquad
 d\leq\lfloor a^4/256\rfloor.
\]
These bounds follow from the trace inequality and Maclaurin's
inequalities. Factoring every polynomial in this finite box and
counting its positive roots exactly gives no irreducible totally
positive quartic. The complete check is included in the verifier.
If $r=2$, a totally positive irreducible quadratic algebraic integer
has trace at least three, so its trace from $K$ is at least six,
contradicting the same bound. Thus $r=1$, and
$4\delta+2\leq B$ forces $\delta=180$.

There would then be four signed Casimir eigenvalues $180$, one
of value two, and two further positive eigenvalues $u,v$.
The reciprocal-codegree identity gives
$90(u+v)=43uv$. Their residual quadratic has integer coefficients,
so it is
\[
 z^2-43kz+90k\qquad(1\leq k\leq7).
\]
Each of these seven quadratics is irreducible, while
$(43k)^2/(90k)=1849k/90$ is not integral. Their roots are therefore
not $d$-numbers, contrary to the formal-codegree theorem
\cite{OstrikCodegrees}; doubling a codegree does not change the
$d$-number property. This excludes $e=4$.

Consequently $e=2$ and $K=\mathbb Q(\sqrt5)$. Write
$\beta=(a+b\sqrt5)/2$, with $a,b\in\mathbb Z$ of the same parity.
Total positivity and the trace bound give
$1\leq a\leq5$ and $a^2>5b^2$. The six possibilities are
$(a,b)=(2,0),(3,\pm1),(4,0),(5,\pm1)$, proving the assertion.
\end{proof}

The exact arithmetic underlying this reduction is preserved in
\begin{quote}\small
\path{evidence/upper_residual_20261002c/scoped_research/QUADRATIC_REMAINDER.md}
\end{quote}
at source commit \texttt{1592905dc573227ad03659a6370a9459d40db26c}.
The present verifier independently replays its complete quartic box,
seven quadratic $d$-number failures, and source Casimir checks.
